\scanpage{63}

On considère $A_M$ comme l'ens des \uline{arêtes} d'une structure de graphe sur $S_M$. \deleted{Le groupe} On a une application naturelle

\[
 A_M \xrightarrow{\unsure{\tau|t}\text{ ou }\tau_A} \deleted{\unsure{\Delta|A}}\Pi_M \overset{\text{déf}}{=} (\text{dét}\,M)^*/{\pm 1_\Lambda}
 \tag*{(161)}
\]

% slanted
\margin{on regardera $\Pi_M$ comme un ens. quotient de $A_M$, \unsure{grâce} à $\tau$.}

déduite par passage au\deleted{\unsure{x}} quotient de l'application

\[
 \left\{\begin{aligned}
  \widetilde A_M &\xrightarrow{\widetilde\tau} \Pi_M \\
  \{x,y\} &\longmapsto x\wedge y \bmod\text{signe}
 \end{aligned}\right.
 \tag*{(162)}
\]

\margin{NB. ne change pas si on échange $x$ et $y$ ou si on change $x$ ou $y$ ou les deux de signe}

Posons, pour $i\in\Pi_M$

\[
 A_{M,i}=\tau^{-1}(\{i\})\ , \tag*{(163)}
\]

qu'on considère comme l'ens des arêtes pour une structure de graphe \deleted{\unsure{sur}} $\Sigma(i)$ sur $S_M$.

Le groupe $\operatorname{Gl}_\Lambda(M)$ opère sur $(\deleted{\unsure{A}}\unsure{S|A}_M,A_M)$ par transport de structure, ainsi que sur $\Pi_M$, et $\tau$ (161) commute aux opérations de $\operatorname{Gl}_{\unsure{\Lambda|\lambda}}(M)$. Comme le sous-groupe $\{\pm\operatorname{id}_M\}$ de $\operatorname{Gl}_\Lambda(M)$ opère trivialement sur les ens $S_M,\Pi_M$, on

% archive sheet number: 62
